% ============================================================
%  04-resume.tex  —  Résumé (FR) + Abstract (EN) — Version Premium
% ============================================================
\chapterstar{Résumé}
\thispagestyle{fancy}

\vspace{6pt}

% ── Résumé français ────────────────────────────────────────
\begin{tcolorbox}[
  enhanced, breakable,
  colback=sectionbg, colframe=isetblue,
  arc=5pt, boxrule=0.8pt,
  attach boxed title to top left={yshift=-2mm, xshift=8mm},
  boxed title style={
    colback=isetblue, colframe=isetblue,
    arc=3pt, boxrule=0pt,
    fontupper=\bfseries\color{white}
  },
  title={Résumé},
  left=10pt, right=10pt, top=10pt, bottom=10pt
]
\setlength{\parskip}{4pt}
Le secteur des énergies renouvelables, et particulièrement le photovoltaïque,
connaît une croissance soutenue en Tunisie dans le cadre du Plan Solaire
Tunisien. Pourtant, les sociétés d'installation de panneaux solaires peinent
encore à industrialiser leurs processus commerciaux et techniques, faute d'un
outil unifié et adapté à leurs besoins métiers.

Le présent Projet de Fin de Parcours présente la conception et la réalisation
de \textbf{SolarEase}, une plateforme web intelligente dédiée aux installateurs
photovoltaïques. La solution centralise l'intégralité du cycle de vie d'un
projet solaire~: authentification sécurisée par \textbf{JWT et OTP}, gestion
des clients et des projets, calcul automatisé de dimensionnement intégrant
l'API \textbf{PVGIS} pour des données d'irradiance réelles, comparaison de
scénarios \textit{Classique versus Night~Panel} avec stockage batterie,
analyse financière sur 25~ans, génération de rapports \textbf{PDF}
professionnels et aide à la décision par \textbf{Intelligence Artificielle}
(LLM Ollama / TinyLLaMA en mode local).

L'architecture adoptée est de type \textbf{microservices} (Spring Boot~3.2 /
Java~17), organisée en quatre services indépendants\,---\,\textit{Identity,
Project, Dimensioning} et \textit{Gateway}\,---\,chacun disposant de sa propre
base de données \textbf{PostgreSQL~16}. Le frontend est développé en
\textbf{React~18 / TypeScript} avec Vite et TailwindCSS~; le déploiement est
orchestré via \textbf{Docker Compose}.

Les résultats valident la faisabilité technique et opérationnelle de la solution~:
temps de traitement réduit, traçabilité accrue et propositions commerciales
de qualité professionnelle.
\end{tcolorbox}

\vspace{4pt}
\noindent\small\textbf{Mots-clés~:}
\textcolor{isetblue}{photovoltaïque} $\cdot$
\textcolor{isetblue}{microservices} $\cdot$
\textcolor{isetblue}{Spring Boot} $\cdot$
\textcolor{isetblue}{React} $\cdot$
\textcolor{isetblue}{dimensionnement solaire} $\cdot$
\textcolor{isetblue}{PVGIS} $\cdot$
\textcolor{isetblue}{Night Panel} $\cdot$
\textcolor{isetblue}{intelligence artificielle} $\cdot$
\textcolor{isetblue}{Docker}

\vspace{16pt}
\begin{center}
  \textcolor{isetgold}{\rule{0.15\textwidth}{0.6pt}}
  \enspace\textcolor{isetblue!40}{\small$\blacklozenge$}\enspace
  \textcolor{isetgold}{\rule{0.15\textwidth}{0.6pt}}
\end{center}
\vspace{12pt}

% ── Abstract anglais ───────────────────────────────────────
\begin{tcolorbox}[
  enhanced, breakable,
  colback=sectionbg, colframe=isetgold,
  arc=5pt, boxrule=0.8pt,
  attach boxed title to top left={yshift=-2mm, xshift=8mm},
  boxed title style={
    colback=isetgold, colframe=isetgold,
    arc=3pt, boxrule=0pt,
    fontupper=\bfseries\color{white}
  },
  title={Abstract},
  left=10pt, right=10pt, top=10pt, bottom=10pt
]
\setlength{\parskip}{4pt}
The renewable energy sector, particularly photovoltaics, is experiencing
sustained growth in Tunisia as part of the national Solar Plan. However, solar
panel installation companies still struggle to industrialize their commercial
and technical processes due to the lack of an integrated tool tailored to their
specific business needs.

This final-year project presents the design and implementation of
\textbf{SolarEase}, an intelligent web platform dedicated to photovoltaic
installers. The solution centralizes the complete lifecycle of a solar
project~: \textbf{JWT and OTP}-based secure authentication, customer and
project management, automated sizing calculations integrating the \textbf{PVGIS}
API for real irradiance data, Classic vs.~Night~Panel scenario comparison with
battery storage, 25-year financial analysis, professional \textbf{PDF} report
generation, and \textbf{AI}-assisted decision support (Ollama / TinyLLaMA LLM
running locally).

The adopted architecture is \textbf{microservices}-based (Spring Boot~3.2 /
Java~17), structured into four independent services\,---\,\textit{Identity,
Project, Dimensioning} and \textit{Gateway}\,---\,each with its own dedicated
\textbf{PostgreSQL~16} database. The frontend is built with
\textbf{React~18 / TypeScript}, Vite and TailwindCSS; deployment is
orchestrated via \textbf{Docker Compose}.

The results validate the technical and operational feasibility of the solution~:
reduced processing time, enhanced traceability, and professional-quality
commercial proposals.
\end{tcolorbox}

\vspace{4pt}
\noindent\small\textbf{Keywords~:}
\textcolor{isetblue}{photovoltaics} $\cdot$
\textcolor{isetblue}{microservices} $\cdot$
\textcolor{isetblue}{Spring Boot} $\cdot$
\textcolor{isetblue}{React} $\cdot$
\textcolor{isetblue}{solar sizing} $\cdot$
\textcolor{isetblue}{PVGIS} $\cdot$
\textcolor{isetblue}{Night Panel} $\cdot$
\textcolor{isetblue}{artificial intelligence} $\cdot$
\textcolor{isetblue}{Docker}

\cleardoublepage
